\Needspace{18\baselineskip}
\section{Counting simulation and its soundness}

\begin{definition}[Up-to closure]
Given a relation $R$ on states and cost orders $P,S$, define
\[
 \upto{P}{S}{R}(s,t)\quad\Longleftrightarrow\quad
 \exists u,v\;\bigl(s\preceq_P u\ \land\ R(u,v)\ \land\ v\preceq_S t\bigr).
\]
\end{definition}
\lean{Simulation}{10}{11}{Machine.UpTo}
The witnesses occur in path order. In particular the right-hand refinement
is from $v$ to $t$, not from $t$ to $v$. Auxiliary refinement can change the
states at which the candidate relation is resumed.

\begin{definition}[Counting simulation]
A relation $R$ is a $(Q,P,S)$-counting simulation if whenever $R(s,t)$:
\begin{enumerate}[label=\textup{(\Alph*)}]
\item If $s$ is final, then $t$ is final.
\item For every transition $s\to s'$, one of the following holds:
\begin{enumerate}[label=\textup{(\roman*)}]
\item There exist $k,z$ such that $s'\steps{k}z$ and $z$ is stuck.
\item There exist $k,l,h_1,h_2$ such that
\[
 s'\steps{k}h_1,\qquad t\steps{l}h_2,\qquad
 (1+k)\,Q\,l,\qquad \upto{P}{S}{R}(h_1,h_2).
\]
\end{enumerate}
\end{enumerate}
\end{definition}
\lean{Simulation}{14}{19}{Machine.CountingSimulation}
The leading $1$ charges for the already prescribed transition $s\to s'$.
Neither $k$ nor $l$ is required to be positive. The stuck alternative is a
deliberate part of the definition, consistent with observing only successful
termination. The finality condition is stronger than simply allowing the
target to reach a final state after some extra steps: it requires the target
to be final already when the source is final.

\begin{theorem}[Soundness of counting simulation]\label{thm:soundness}
Suppose $M$ is deterministic, $(Q,P,S)$ is reasonable, and $R$ is a
$(Q,P,S)$-counting simulation. Then
\[
 R(s,t)\ \Longrightarrow\ s\preceq_Q t.
\]
\end{theorem}
\lean{Simulation}{24}{51}{Machine.CountingSimulation.sound}
This is the abstract transition-system form of the paper's Proposition~37.
Its proof is a strong induction on the length of a successful source
execution. Although simulation provides a step-wise, coinductive style of
reasoning, the soundness theorem itself is proved by induction on a finite
observation; no coinductive datatype is introduced in this file.

\begin{proof}
We prove, simultaneously for all states $s,t,z$, that if $R(s,t)$,
$s\steps{N}z$, and $\fin(z)$, then there are $m,w$ such that
$N\,Q\,m$, $t\steps{m}w$, and $\fin(w)$. Use strong induction on $N$;
the induction hypothesis applies to every smaller natural number and every
pair of related states.

\emph{Base case.} If $N=0$, then $s=z$ is final. Condition~(A) makes
$t$ final. Choose the zero-step execution at $t$ and use $0\,Q\,0$.

\emph{Positive case.} Write $N=n+1$ and decompose the successful execution:
\[
 s\to s'\steps{n}z.
\]
Apply condition~(B) of simulation. Its stuck alternative would give a
finite execution from $s'$ to a stuck state. This is impossible by the
lemma that a state on a successful execution cannot be stuck, applied to
the successful execution $s'\steps{n}z$.

Consequently there are $k,l,h_1,h_2,u,v$ satisfying
\begin{gather*}
 s'\steps{k}h_1,\qquad t\steps{l}h_2,\qquad (1+k)\,Q\,l,\\
 h_1\preceq_P u,\qquad R(u,v),\qquad v\preceq_S h_2.
\end{gather*}
The prefix lemma, using determinism and finality of $z$, gives a remainder
length $r$ such that
\[
 n=k+r,\qquad h_1\steps{r}z.
\]
Apply $h_1\preceq_P u$ to this remaining successful execution. It supplies
$a,z_a$ with
\[
 r\,P\,a,\qquad u\steps{a}z_a,\qquad \fin(z_a).
\]
The decrease condition of reasonableness gives $a\leq r$. Since
$N=1+k+r$, we have $a<N$. This is the strict inequality needed to invoke
the strong induction hypothesis. The proof would not be justified for
arbitrary auxiliary cost preorders.

Applying that hypothesis to $R(u,v)$ gives $b,z_b$ such that
\[
 a\,Q\,b,\qquad v\steps{b}z_b,\qquad \fin(z_b).
\]
Now apply $v\preceq_S h_2$. We obtain $c,z_c$ with
\[
 b\,S\,c,\qquad h_2\steps{c}z_c,\qquad \fin(z_c).
\]
The sandwich condition combines $r\,P\,a$, $a\,Q\,b$, and $b\,S\,c$
to yield $r\,Q\,c$. Additivity of $Q$ then gives
\[
 (1+k+r)\,Q\,(l+c),\quad\text{that is,}\quad N\,Q\,(l+c).
\]
Concatenate $t\steps{l}h_2$ with $h_2\steps{c}z_c$. The result is a
successful execution of length $l+c$ from $t$, which supplies the required
witnesses.
\end{proof}

The code follows these witnesses exactly: its \texttt{count} is $n=k+r$,
\texttt{smaller} is $a<n+1$, \texttt{rc} is $r\,Q\,c$, and
\texttt{total} is $(n+1)\,Q\,(l+c)$. The inclusion $P\subseteq S$ is
not needed in this particular proof; it is part of the shared reasonable
triple and will be used by the specimen argument.
